\section{Introduction}
\label{sec:intro}

Dense per-voxel annotation is the binding constraint on CT segmentation. A
fully supervised model such as TotalSegmentator already segments more than a
hundred anatomical structures at a level of agreement comparable to an expert
reader \cite{wasserthal2023totalsegmentator}, so accuracy is not the open
problem; that performance instead rests entirely on expert voxel-level labels
for every structure of interest. Because it is the labels rather than the
architecture that determine which structures such a system can recognize,
extending it to a new structure, a new protocol, or a new clinical question
means commissioning new annotation.

That commissioning is expensive because voxel-level delineation of CT is not
unskilled work. It requires a reader trained in cross-sectional anatomy who can
identify organ boundaries on axial slices where soft-tissue contrast is often
poor and where adjacent structures share intensity, and such expertise is
correspondingly scarce and costly to secure. The burden also scales
multiplicatively rather than additively, since every structure must be traced on
every slice of every volume, and a single CT volume comprises hundreds of
slices. As a result, densely labeled CT datasets are few in number, and each
one is fixed in the label set that a particular funding decision chose to
support \cite{ji2022amos, luo2022word, rister2020ctorg}, which means a structure
outside that set is simply unavailable to anyone working from it. A second
consequence follows from the same difficulty: where expert judgment is genuinely
ambiguous, as it is at low-contrast boundaries, even careful annotators
disagree, and therefore the ground truth against which these models are trained
and scored itself carries reader variability \cite{nikolov2021nsd}. The
limitation that matters here is consequently not model capacity but the
annotation itself, which raises the question of whether the supervision might
instead be drawn from a source that already exists rather than one that has to
be created.

Text is the natural candidate for such a source, because descriptions of
anatomy are already abundant in clinical reports and in the biomedical
literature, and vision--language pretraining offers a means of learning from
them. Models in the CLIP family learn visual concepts from natural-language
supervision alone \cite{radford2021clip, zhai2023siglip}, and biomedical
adaptations such as PMC-CLIP \cite{lin2023pmcclip} and BiomedCLIP
\cite{zhang2025biomedclip} bring that capability into the medical domain at
scale. If a text prompt could select a region as well as it selects a label,
segmentation would need no pixels at all.

The obstacle is that these models are trained for global alignment. The
contrastive objective relates a pooled image embedding to a pooled caption
embedding, and nothing in it requires an individual patch token to be spatially
discriminative. The conventional way to recover a spatial map from a trained
network is the CAM family \cite{zhou2016cam, selvaraju2017gradcam}, which
extracts one post hoc from a classifier; that route is closed here, however,
because a contrastively pretrained model makes no per-organ decision for a CAM
to explain. What remains available is to backpropagate the image--text
similarity to the input pixels, and that gradient saliency recovers almost no
localization: applied to the pretrained model, it reaches 0.003 mIoU.

This work makes the spatial property an explicit training objective instead.
That objective aligns each of the 196 ViT patch tokens with the organ's text
embedding against the downsampled ground-truth mask, so that per-patch cosine
similarity becomes directly readable as a segmentation map at inference. Because
the map is read straight off the forward pass, the method requires no gradient
computation, no classifier, and no second-stage mask refiner. The contributions
are as follows.

\textbf{A patch-level alignment objective that makes text-prompted localization
work.} Aligning each patch token directly with the organ's text embedding
reaches 0.653 mIoU on 11 organ classes from TotalSegmentator while requiring no
per-pixel labels at inference, and the patch loss contributes more of that
localization than any other component of the method
(\S\ref{sec:results:phase1}).

\textbf{A prompt ablation that establishes what the text is actually doing, and
what it is not.} Corrupting the positional content of the prompts collapses
precisely the classes that depend on laterality while leaving
uniquely-positioned organs untouched, which establishes that the prompt is
load-bearing rather than acting as a class index
(\S\ref{sec:results:ablation}).

\textbf{A safety-relevant failure mode.} An accidental two-word error in one of
eleven prompts produced not a degraded output but a confident, clean
segmentation of the wrong organ, which means that failures of this method are
directional rather than degradative and are therefore unsignaled by the output
itself (\S\ref{sec:results:ablation}).

\textbf{A downstream clinical application of the frozen model.} Carried
unmodified onto Spine-Mets-CT-SEG \cite{pieper2024spinemets}, a dataset and
acquisition protocol never seen in training, the same weights rank the organ
nearest a vertebral metastasis well above chance across 180 lesions in 46
patients, at $p < 0.0001$ under a patient-clustered permutation test
(\S\ref{sec:results:lesion}).

\textbf{Negative results with diagnosed mechanisms.} Neither of two
pre-registered calibration variants was adopted, in each case for an identified
reason rather than for an unexplained regression
(\S\ref{sec:results:negatives}) -- a diagnosis that extends to the residual
localization failure itself, which is not diffuse but concentrates in two of
eleven classes (\S\ref{sec:discussion}).

We also report a ground-truth error that we found in our own evaluation table
and corrected from anatomical literature, together with a calibration leak that
we found and repaired. Because both of these changed headline numbers,
reporting them is part of what the paper claims rather than an aside to it.
